\section{The completions and Hodge jumps at the ramified primes}

% BEGIN R79Q1
\paragraph{R79Q1. Inertia and the CM type at the two primes.}
Let $D$ be an original datum, $E$ its host field and
$\ell\in\{q,d\}$. Both primes are $3$ modulo $4$ by the host
congruences. For every place $w\mid\ell$ the ramification index is
two and complex conjugation takes $w$ to a different place.
At $q$ every element of the decomposition group preserves the CM
type on embeddings. At $d$ every nonidentity element of inertia
interchanges the CM type and its conjugate.

Indeed the host identifies $q$-inertia with the Artin image of
units congruent to one modulo $Q_Sd$ in conductor $Q_S d q$.
This kernel embeds in the cyclic group of units modulo $q$ by
the Chinese remainder theorem. Its image has exponent at most two
by the host's $q$-character condition. The unit equal to $-1$
modulo $q$ and to one modulo $Q_S d$ has first quadratic sign
$(-1/q)=-1$. Hence its Artin image has order two and generates
inertia. At $d$ the host's inertia identity gives exactly
$\{1,\iota_d\}$, with $\iota_d$ acting by minus one on
$\sqrt{-d}$. In a Galois extension the order of inertia at a
finite place equals its ramification index, giving the assertion at
both primes. Both divide the conductor, where the host's
decomposition condition excludes complex conjugation from the
stabilizer. At $q$ the host's decomposition group belongs to
$\Psi$ and fixes $\sqrt{-d}$. An embedding belongs to the chosen
CM type precisely when it sends this generator to $+i\sqrt d$.
The other possible value is $-i\sqrt d$, distinct since $d>0$.
This proves both assertions about the CM type, hence all of the
original $\mathrm{Conclusion1}(D)$.
% END R79Q1

% BEGIN R79Q2
\paragraph{R79Q2. The completion inside the first cyclotomic layer.}
Let $K$ be the local field in the given bridge at $w\mid\ell$,
let $K_0$ be its maximal unramified subfield, and put
$K_1=K_0(\mu_\ell)$. The image $y$ of the host generator
$\sqrt{-q}$ or $\sqrt{-d}$, as appropriate, satisfies $y^2=-\ell$.
The bridge embeds $E$ densely in $K$. The restricted $\ell$-adic
absolute value on $E$ induces a place above $\ell$, so $K$ is
its completion. R79Q1 applies to every such place.
For finite extensions of $\mathbb Q_\ell$ the residue extension
lifts to the maximal unramified subfield, so $[K:K_0]$ equals the
ramification index, two by R79Q1. The normalized valuation of $y$
is $1/2$, whereas the valuation group of $K_0$ is $\mathbb Z$.
Thus $K=K_0(y)$.

To see the inclusion in the specified $K_1$, take a primitive
$\ell$th root $\zeta$ and the quadratic character $\lambda$ of
$\mathbb F_\ell^\times$, extended by zero. For
$G=\sum_a\lambda(a)\zeta^a$, replacing $\zeta$ by $\zeta^{-1}$
multiplies $G$ by $\lambda(-1)=-1$. The product of the two sums
is $\ell$: write $a=ub$ with $b\ne0$ in the double sum. The
inner sum $\sum_{b\ne0}\zeta^{(u-1)b}$ is $\ell-1$ for $u=1$
and $-1$ otherwise, while $\sum_u\lambda(u)=0$. Hence
$G^2=-\ell$ and $y=\pm G$ belongs to $K_1$. The polynomial
$\Phi_\ell(1+X)$ is Eisenstein over the unramified $K_0$.
It follows that $[K_1:K_0]=\ell-1$ and
\[
 K\subseteq K_1,\qquad
 [K_1:\mathbb Q_\ell]=\frac{\ell-1}{2}[K:\mathbb Q_\ell].
\]
These are the four assertions of $\mathrm{Conclusion0}(D)$ on
the original bridge.
% END R79Q2

% BEGIN R79Q3
\paragraph{R79Q3. The signed Hodge windows over every coefficient extension.}
Use the member, hub $X$, good family $R=X.r$ and
$\mathrm{AdmJ}(L)$ from F78D2. Take the single sign parameter
$e$ in its archimedean dictionary. Fix any twisted $i$ at $q$ or $d$
satisfying the native nonemptiness condition $\mathrm{HG}$, any
$w\mid\ell_i$, any coefficient extension $C$, and any embedding
$\tau:K\to C$ over $\mathbb Q_{\ell_i}$. Composing the host
embedding at $w$ with $\tau$ and the given coefficient identification
with $\mathbb C$ gives the label $g$ and its archimedean place in
$\mathrm{AdmRead}$. Call the sign positive when
$(g\in\Psi)\Leftrightarrow(e=\mathrm{true})$.

Put $T=(t-1)/2$, $v=e_{j-1}$ with $j=j_i$, $\eta=L.\mathrm{etaSz}$,
$S=L.\mathrm{Sc}$, and $b=L.\mathrm{sc}.b$. The type primes
split completely, so R79Q1 excludes $q,d$ from them. The
away-from-type de Rham reading and the de Rham/Weil--Deligne
identity in $\mathrm{AdmJ}$ therefore apply at every point $x$.
The given coefficient extension is the one allowed in
$\mathrm{AdmRead}$. As in the ambient-weight calculation of
S79P3, subtract the weight of $\chi_1$ from the weights at that
embedding. For positive sign this gives the three lists
\[
 T+v+v_a\ (0\le a<k-1),\qquad
 T+v-b_a\ (0\le a<m),\qquad 0,\ldots,t-1.
\]
For negative sign apply $z\mapsto2T-z$ to all entries.
The dictionary's criterion is nonzero multiplicity in this
shifted multiset. In particular it computes the jumps of
$R_x\chi_1^{-1}$ at every $x$, including outside the good locus.
All entries are integral because $t$ is odd and the source
parameters $v,v_a,b_a$ are integral. Write $J_c,J_l$ for the
sets from the first two lists in the chosen orientation.
The $v_a$ are strictly decreasing and the $b_a$ strictly
increasing by the archimedean parameter inequalities. Consequently
$|J_c|=k-1$, $|J_l|=m$.

For positive sign take $s=v+T$, and for negative sign take $s=v-T$.
The family windows put $v$ between $w''$ and $b+\eta+1$
by G79E3. Their lower bound $w''\ge m(t+1)\eta$, together
with $m\ge2$, $t>3m$ and $\eta\ge t$, gives
$w''>\eta+t+1+T$. Thus $\eta+t+1<s\le b+\eta+t+1$.
Since $0\le v_a\le S$, the signed cloud lies in $[s,s+S]$ or
$[-s-S,-s]$ respectively. The low differences have
$T+2=(t+3)/2\le |v-b_a|\le\eta$ by S79P4, so in either
orientation
\[
 J_l\subseteq[T-\eta,-2]\ \cup\ [t+1,T+\eta].
\]
These intervals are exchanged by reflection, and the integer
character string is preserved. This proves every exclusion at
$-1,t$ and in the character string for a low entry.

Finally, if the source's $\mathrm{UpAt}(e,C,\tau)$ holds, it gives
a base-coefficient embedding $\tau_0$ inducing $\tau$ under
$C$'s inclusion; $\mathrm{DownAt}$ gives the opposite base type.
The bridge and the original coefficient map agree on the
coefficient DVR, hence on its fraction field. Evaluating their common
complex embedding on the host generator $\sqrt{-d}$ therefore
reads the same sign for $\tau_0$ and for $g$. A CM-type base
embedding gives $+i\sqrt d$ and its conjugate gives $-i\sqrt d$.
Consequently $\mathrm{UpAt}$ forces the positive alternative
for $e$ and $\mathrm{DownAt}$ the negative one. Even when neither
antecedent is given, $g$ supplies one orientation. With that
orientation the jump set at every $x$ is precisely
$J_c\cup J_l\cup\{0,\ldots,t-1\}$. This proves the original
distinguished-prime Hodge window for the single chosen $e$.
% END R79Q3

% BEGIN R79Q4
\paragraph{R79Q4. The degree and declaration bounds for the ramified count.}
At either prime declare all $w\mid\ell$ to be rows and no places
to be columns. This is a finite declaration covering exactly the
required places. For any subset $A$ of these places,
\[
 2|A|\le\sum_{w\in A}[E_w:\mathbb Q_\ell]\le
 [E:\mathbb Q]=2f .
\]
Indeed each degree is twice its positive residue degree by R79Q1,
and the sum over all places is the global degree.
R79Q2 gives $[E_w(\mu_\ell):E_w]=(\ell-1)/2$. The global degree
$[E(\mu_\ell):E]$ is at least this local degree by completing at
a place above $w$. The host size condition $4tr+1<\ell$ thus
gives $tr<[E(\mu_\ell):E]$. Finally $\ell$ is odd and outside
the host's small set, so the native local level is one. As
$1\le j<t<\ell$, every $v_\ell(j)$ is zero. The defining sum
for $v_\delta(i)$ is then $t-1$, in particular $v_\delta(i)\le t$.

R79Q1--R79Q3 supply the local fields and Hodge windows used in
the distinguished-prime count. The $\mathrm{Conclusion2}$ required
by J78H2 still calls for the integral local lifting and character
block bounds on these rows, and their global lattice-index
inequality for the same representations. In particular its source
uses the local block witness
$c_B\le60tf(r+t)\log b$ at each allowed point, together with
the row-lifting estimate with constant $160$. The field-degree
and jump identities alone supply neither length estimate.
% END R79Q4
